\section{Tercera iteración de la Fase 2: El estado vive en la lista de segmentos}
\label{sec:EstadoEnLista}

La iteración anterior (F2-I2) dejó al motor de juego con una secuencia fija que solo se
anunciaba por consola: el estado del juego era un número suelto. Esta iteración le da su
forma definitiva de trabajo: \textbf{el motor de juego ESCRIBE el estado en la lista de
	segmentos en memoria}, la misma que la carga del GWI dejó preparada. Es el paso clave del
contrato de la Fase 2 ---«la lista ES el estado»--- hecho carne: a partir de aquí, todo lo
que el juego sepa sobre sí mismo vive en los objetos de esa lista, y quien quiera conocer
el estado (hoy la consola, mañana el motor gráfico, algún día el ESP32) solo tiene que
leerla. Nada de estados paralelos ni copias: una lista, un escritor, muchos lectores.

\subsection{Qué se construye}
\label{subsec:QueSeConstruyeF2I3}

\begin{itemize}
	\item \textbf{Las dos banderas por segmento.} Cada segmento de la lista en memoria
	gana dos banderas, las mismas que definió el análisis de F2-I1:
	\texttt{visible} (el segmento está encendido en el cristal LCD) y
	\texttt{cambiar} (el segmento cambió desde la última presentación: la ``suciedad''
	que el motor gráfico usará para saber qué repintar). Las banderas no se añaden al
	GWI: son \emph{estado de ejecución}, nacen y mueren con la partida; el fichero sigue
	siendo estático.
	\item \textbf{El tick del motor de juego, con solo dos acciones.} En cada tick el
	motor de juego hace exactamente dos cosas, en este orden:
	\begin{enumerate}
		\item \textbf{Apaga las banderas de cambiado}: todo segmento con
		\texttt{cambiar} a \texttt{true} pasa a \textbf{false}. Con esto queda
		declarado que el estado anterior ya fue presentado (cuando exista el motor
		gráfico, esta es la señal de ``me han leído'').
		\item \textbf{Cambia \texttt{visible} en el segmento de la secuencia}: avanza
		la secuencia fija, apaga el segmento que estaba encendido
		(\texttt{visible=false}, \texttt{cambiar=true}) y enciende el nuevo
		(\texttt{visible=true}, \texttt{cambiar=true}).
	\end{enumerate}
	Nada más: sin teclado, sin reglas, sin PRNG. El resto de la lista no se toca.
	\item \textbf{Depuración con llave: \texttt{GW\_DEBUG}.} Nace una variable global
	\texttt{GW\_DEBUG} en el \texttt{index.html}, editable a mano, que gobierna todos
	los mensajes de diagnóstico de los motores: con ella activa, el motor de juego
	cuenta en consola las dos cosas que hace (qué banderas de cambiado apaga, qué
	segmento enciende o apaga); con ella apagada, la consola queda limpia y el
	sistema corre en silencio. Los mensajes son \emph{diagnóstico del sistema}, no
	salida del juego: la misma llave servirá en las iteraciones siguientes para
	cualquier motor que quiera contar lo que hace.
\end{itemize}

\subsection{Detalles de realización}
\label{subsec:DetallesF2I3}

\begin{itemize}
	\item \textbf{La lista en memoria, expuesta.} La lista de segmentos que la carga del
	GWI deja en memoria vive (bien encapsulada) en el módulo de autoría
	(\texttt{BaseDescomponer}); esta iteración la expone con un único accesor de
	lectura que devuelve \emph{la referencia} al array ---no una copia---, porque el
	contrato es escribir en ella. El motor de juego la pide al arrancar; si no hay
	cartucho GWI cargado la lista está vacía y el motor no revienta (avisa si hay
	depuración).
	
	\item \textbf{El motor de juego sigue siendo el único escritor}. Ni el director ni
	el motor gráfico tocan las banderas; el gráfico (todavía de anuncio en consola)
	sigue siendo lector puro. La disciplina de un solo escritor se mantiene desde el
	primer día para que nunca haya que depurar dos escritores.
	
	\item \textbf{Arrancar limpia el estado}: al lanzar el juego, todas las banderas de
	la lista se ponen a \texttt{false} y la secuencia empieza desde el principio; el
	reloj y el estado se reinician juntos, como quedó en F2-I2.
	
	\item \textbf{Verificación}: la simulación comprueba que las banderas cambian de
	verdadero en los objetos de la lista (el segmento encendido lleva
	\texttt{visible=true, cambiar=true}, el anterior \texttt{visible=false,
		cambiar=true}, el resto sin tocar), que el apagado de \texttt{cambiar} ocurre al
	tick siguiente, que con \texttt{GW\_DEBUG} apagada no aparece mensaje alguno, y
	que el bucle largo es estable.
\end{itemize}

\subsection{Ejecución}

Cargamos un setup, luego un cartucho, y al pulsar sobre START vemos que el contador de ticks
aumenta y aparecen los mensajes en la consola..

\begin{figure}[H]
	\centering
	\includegraphics[width=0.9\linewidth]{imagenes/propias/C05-I020}
	\caption[]{Ejecución del motor de gráfico básico.}
	\label{fig:c05-i020}
\end{figure}


\clearpage


